\scanpage{171}
\mspage{87}

\unsure{qui|Qui} se réconstitue à partir du diagramme de groupes partiel

\diagram{\begin{tikzcd}[ampersand replacement=\&,row sep=small]
(44)\quad\&\widetilde B''\&\\
\&\widetilde U\arrow[u,hook,"Z"]\arrow[r,hook,"Z'"]\&B'
\end{tikzcd}}

sauf que pour en tirer l'hom horizontal $i$ de (43), dont la composante suivant $Z$ doit être l'hom can de passage au quotient, il faut se donner la première composante, savoir un homomorphisme de groupes

\[ \widetilde B''\longrightarrow B' \]

la commutativité de (43) signifiant que ce dernier rend commutatif le triangle

\diagram{\begin{tikzcd}[ampersand replacement=\&,row sep=small]
(45)\quad\&\widetilde B''\arrow[dr]\&\\
\&\widetilde U\arrow[u,hook,"Z"]\arrow[r,hook,"Z'"]\&B'
\end{tikzcd}}

de telle sorte qu'elle donne \inserted{donc} naissance à un hom. d'extensions de groupes

\diagram{\begin{tikzcd}[ampersand replacement=\&,column sep=small,row sep=small]
(\unsure{45|46})\quad 1\arrow[r,dash]\&\widetilde U\arrow[r]\arrow[d,"1"]\&\widetilde B''\arrow[r]\arrow[d,"\psi"]\&Z\arrow[r,dash]\arrow[d]\&1\\
1\arrow[r,dash]\&\widetilde U\arrow[r]\&B'\arrow[r]\&Z'\arrow[r,dash]\&1.
\end{tikzcd}}

Ainsi, les données (42 a)b)) se réduisent à la donnée d'un hom de suites exactes (46), ou encore d'un diagramme d'hom de groupes

\diagram{\begin{tikzcd}[ampersand replacement=\&,row sep=tiny]
(47)\quad\&\&Z\arrow[d]\\
\&B'\arrow[r,"\text{épi}"']\&Z',
\end{tikzcd}}

% archive sheet number: 170
